\section{Conclusion}\label{sec:conclusion}

This study transforms a GA-based conference precursor into a controlled benchmark for neural intraday trend classification. It preserves the predecessor's practical insight---a compact, Information-Gain-selected technical representation can support metaheuristic hyperparameter search---but replaces randomized validation and a single-optimizer design with two primary chronological holdout protocols, a separate exploratory temporal cross-validation analysis, five budget-matched optimizers, and locked predictive and economic assessment. The five-minute feature vector at bar $t$ is paired with the regime label at $t+1$, derived from the 15-minute peak-and-trough regime sequence.

The current evidence indicates configuration-dependent changes between the two primary fitness protocols, which differ in both metric and train/validation composition. MLP leads the predictive aggregates under weighted-accuracy fitness, while the aligned next-bar replay gives 1D-CNN the largest mean net profit under both protocols; RF has the least negative mean drawdown under MCC/$F_1$ fitness. These results reinforce the central methodological lesson: selection objective, reported metric, and economic utility must be treated as related but distinct quantities.

The manuscript deliberately avoids a confirmatory claim about a universal best fitness function because only three paired seeds are currently available across every cell. The 15-minute pivot labels are retrospective, and the current results do not purge samples whose label-defining regime crosses a temporal split boundary. A purged rerun with preserved pivot-confirmation times and more seeds across multiple market periods is needed before stronger generalization claims. The present benchmark is therefore exploratory evidence for this fixed chronological interval.
